\documentclass[11pt]{article}
\usepackage[margin=1in]{geometry}
\usepackage{amsmath,amssymb,amsthm}
\usepackage{booktabs}
\usepackage{hyperref}
\newtheorem{theorem}{Theorem}
\newtheorem{lemma}[theorem]{Lemma}
\newtheorem{corollary}[theorem]{Corollary}
\newtheorem{proposition}[theorem]{Proposition}
\theoremstyle{definition}
\newtheorem{definition}[theorem]{Definition}
\theoremstyle{remark}
\newtheorem{remark}[theorem]{Remark}
\newcommand{\F}{\mathbb{F}_{11}}
\newcommand{\PG}{\mathrm{PG}(2,11)}
\newcommand{\PGL}{\mathrm{PGL}(3,11)}

\title{Proof of Conjecture 00000001091:\\
the complete $10$-arcs of $\mathrm{PG}(2,11)$ form a single class}
\author{jilint777}
\date{2026-10-04}

\begin{document}
\sloppy
\maketitle

\begin{abstract}
Conjecture 00000001091 asserts that the complete $10$-arcs of the projective
plane $\PG$ form a single isomorphism class. We prove this. Every $10$-arc
contains a frame, and $\PGL$ is regular on ordered frames, so it suffices to
classify the complete $10$-arcs through the standard frame
$e_1,e_2,e_3,(1,1,1)$. The $24$ projectivities permuting the frame split the
$72$ admissible further points into four orbits. A short exhaustive search
from one representative of each orbit finds $13$ complete $10$-arcs, and each
is mapped onto a fixed representative by an explicit matrix. The whole
argument (linear algebra over $\F$, the frame lemma, the symmetry reduction,
the search with its soundness proof, and the $13$ certificates) is formalized
in Lean~4 (core library only, no \texttt{sorry}, no \texttt{native\_decide};
axioms \texttt{propext} and \texttt{Quot.sound} only). An independent Python
check finds $84$ complete $10$-arcs through the frame, all in one orbit, with
stabilizer $A_5$ of order $60$. This agrees with the literature: Ball and
Lavrauw state that $\PG$ has a unique complete arc of size $q-1=10$.
\end{abstract}

\section{The conjecture and how we read it}
\begin{quote}
\emph{Definition:} regular completions of arcs.
\emph{Conjecture:} the complete $10$-arcs of $\PG$ form a single isomorphism
class (an entry of the $\PG$ classification, computationally verified).
\end{quote}
The Chinese version says the same.

A \emph{$k$-arc} of a projective plane is a set of $k$ points, no three of
which are collinear. It is \emph{complete} if it is not contained in a
$(k+1)$-arc, i.e.\ every point outside the arc lies on a secant (a line joining
two points of the arc). The definition line (``regular completions'') only
names the topic; the claim to be decided is the conjecture.

Two arcs are \emph{isomorphic} (equivalent) if a collineation of the plane maps
one onto the other. The collineation group of $\mathrm{PG}(2,q)$ is
$\mathrm{P\Gamma L}(3,q)$ (fundamental theorem of projective geometry). For
$q=11$ prime, the field $\F$ has no nontrivial automorphism, so
$\mathrm{P\Gamma L}(3,11)=\PGL$. Hence ``isomorphic'' means
\emph{projectively equivalent}: some invertible $3\times3$ matrix over $\F$
maps one arc onto the other. This is the notion used throughout, and it is the
strongest one (a single matrix), so the result also holds for any coarser
notion. We prove:

\begin{theorem}\label{thm:main}
Complete $10$-arcs exist in $\PG$, and any two of them are projectively
equivalent. So the complete $10$-arcs form exactly one class.
\end{theorem}

\section{The model}\label{sec:model}
\paragraph{Points and lines.}
A point is a $1$-dimensional subspace of $\F^3$, written as its normalized
vector (first nonzero coordinate $1$). We number the $133$ points:
\[
11a+b\mapsto(1,a,b)\ (0\le a,b\le10),\qquad 121+b\mapsto(0,1,b),\qquad
132\mapsto(0,0,1).
\]
Lines are numbered the same way by their normalized dual coordinates, and point
$p$ lies on line $L$ iff $\mathrm{coord}(p)\cdot\mathrm{coord}(L)=0$. Three
points are \emph{collinear} if they lie on a common line. The standard frame is
\[
\mathcal F=\{e_1,e_2,e_3,u\}=\{0,121,132,12\},\qquad u=(1,1,1).
\]

\paragraph{Action of matrices.}
For an invertible matrix $M$ (with inverse $N$), $M$ maps the point $p$ to the
point of $M\,\mathrm{coord}(p)$, and the line $L$ to the line of
$N^{\mathsf T}\mathrm{coord}(L)$. The identity
\[
(N^{\mathsf T}\ell)\cdot(Mx)=\ell\cdot(NMx)=\ell\cdot x
\]
shows that incidence is preserved.

\paragraph{Representative.} Our representative complete $10$-arc is
\[
R=\{0,121,132,12,25,35,49,74,103,119\},
\]
i.e.\ the points $(1,0,0)$, $(0,1,0)$, $(0,0,1)$, $(1,1,1)$, $(1,2,3)$,
$(1,3,2)$, $(1,4,5)$, $(1,6,8)$, $(1,9,4)$ and $(1,10,9)$.

\section{Proof of Theorem~\ref{thm:main}}

\begin{lemma}[projectivities preserve complete arcs]\label{lem:preserve}
Let $M$ be invertible. Then $M$ permutes the points and the lines and
preserves incidence. In particular, $p,q,r$ are collinear iff $Mp,Mq,Mr$ are.
Consequently, $M$ maps complete $k$-arcs onto complete $k$-arcs.
\end{lemma}
\begin{proof}
The first statements follow from the identity above and $NM=I$. Let $A$ be a
complete $k$-arc. Then $MA$ has $k$ points, since $M$ is injective on points.
It is an arc, since collinearity of images implies collinearity of the
preimages. For completeness, let $p'\notin MA$ and put $p=Np'$. Then $p\notin
A$, so $p$ lies on a secant $qr$ of $A$, and $p'=Mp$ lies on the secant
$(Mq)(Mr)$ of $MA$.
\end{proof}

\begin{lemma}[frame lemma]\label{lem:frame}
Let $a_0,a_1,a_2,a_3$ be four points, no three collinear. Then some invertible
$M$ maps them to $e_1,e_2,e_3,u$ respectively.
\end{lemma}
\begin{proof}
Classically, write $\mathrm{coord}(a_3)=\sum\lambda_i\,\mathrm{coord}(a_i)$.
All $\lambda_i\neq0$ by general position. Then the matrix with columns
$\lambda_i\,\mathrm{coord}(a_i)$ maps $e_1,e_2,e_3,u$ to $a_0,a_1,a_2,a_3$, and
$M$ is its inverse.

In Lean we use an equivalent four-step construction. Each step is a matrix from
an explicit table that was checked exhaustively:
\begin{enumerate}
\item $t_1(p)$ maps $p\mapsto e_1$;
\item $t_2(p)$ fixes $e_1$ and maps $p\mapsto e_2$, for $p\neq e_1$;
\item $t_3(p)$ fixes $e_1,e_2$ and maps $p\mapsto e_3$, for $p\notin e_1e_2$;
\item $t_4(p)$ fixes $e_1,e_2,e_3$ and maps $p\mapsto u$, for $p$ off the
triangle.
\end{enumerate}
By Lemma~\ref{lem:preserve}, the image of $a_{k}$ after $k$ steps satisfies the
hypothesis of step $k+1$.
\end{proof}

Every $10$-arc contains four points, no three collinear. By
Lemmas~\ref{lem:preserve} and~\ref{lem:frame}, every complete $10$-arc is
therefore equivalent to a complete $10$-arc containing $\mathcal F$. It
remains to classify those.

\paragraph{Candidates and the frame stabilizer.}
A point of an arc $S\supseteq\mathcal F$ outside $\mathcal F$ lies on none of
the six sides of the quadrangle $\mathcal F$. There are exactly $72$ such
points. The $24$ projectivities $g_\sigma$ mapping $(e_1,e_2,e_3,u)$ to a
permutation of itself form a group $G_{\mathcal F}\cong S_4$. It acts on the
$72$ candidates with four orbits:
\begin{center}\small
\begin{tabular}{@{}llp{0.62\textwidth}@{}}
\toprule
orbit & rep. & points\\
\midrule
$O_1$ (12) & $25=(1,2,3)$ & 25, 32, 35, 52, 61, 71, 73, 83, 92, 109, 112, 119\\
$O_2$ (24) & $26=(1,2,4)$ & 26, 28, 29, 31, 39, 43, 46, 51, 63, 65, 68, 69, 75, 76,
79, 81, 93, 98, 101, 105, 113, 115, 116, 118\\
$O_3$ (24) & $27=(1,2,5)$ & 27, 30, 38, 40, 41, 50, 53, 54, 57, 58, 64, 70, 74, 80,
86, 87, 90, 91, 94, 103, 104, 106, 114, 117\\
$O_4$ (12) & $37=(1,3,4)$ & 37, 42, 47, 49, 59, 62, 82, 85, 95, 97, 102, 107\\
\bottomrule
\end{tabular}
\end{center}

\begin{lemma}[symmetry reduction]\label{lem:reduce}
Let $S$ be a complete $10$-arc containing $\mathcal F$. Let $j$ be the least
index with $S\cap O_j\neq\emptyset$, and let $r_j$ be the representative of
$O_j$. Then some $g\in G_{\mathcal F}$ maps $S$ to a complete $10$-arc $S'$
with
\[
S'\supseteq\mathcal F\cup\{r_j\}\quad\text{and}\quad
S'\cap(O_1\cup\dots\cup O_{j-1})=\emptyset .
\]
\end{lemma}
\begin{proof}
Such a $j$ exists, since $S\setminus\mathcal F\neq\emptyset$ consists of
candidates. Pick $c\in S\cap O_j$ and $g$ with $gc=r_j$. Then $g\mathcal F=
\mathcal F$, and the orbits are $g$-invariant. If $x\in S'\cap O_i$ with $i<j$,
then $g^{-1}x\in S\cap O_i$, contradicting the choice of $j$.
Lemma~\ref{lem:preserve} shows that $S'$ is a complete $10$-arc.
\end{proof}

\begin{lemma}[exhaustive search]\label{lem:search}
For each $j\in\{1,2,3,4\}$, let $S$ be a complete $10$-arc with
$S\supseteq\mathcal F\cup\{r_j\}$ and $S\cap(O_1\cup\dots\cup O_{j-1})=
\emptyset$. Then, as a set, $S$ is one of the $13$ arcs listed in
Table~\ref{tab:arcs}.
\end{lemma}
\begin{proof}
Let $C_j=\mathcal F\cup\{r_j\}$. The other five points of $S$ are not in $C_j$,
not in $O_1\cup\dots\cup O_{j-1}$, and on no secant of $C_j$. These points form
an explicit list $\mathrm{rest}_j$ of $43$, $37$, $22$ and $7$ points
respectively.

The search chooses five of them in list order. After choosing $c$ it discards
every later point on a line through $c$ and a chosen point. It also prunes a
branch when too few points remain. At a leaf (ten points) it either finds a
point on no secant, so the leaf is not complete, or it finds the leaf in
Table~\ref{tab:arcs}.

The search is sound. At each node:
\begin{itemize}
\item the chosen points belong to $S$;
\item the remaining list contains every point of $S$ that is not chosen;
\item the number of points of $S$ in the remaining list equals the number still
to be chosen.
\end{itemize}
Branching at the first point of $S$ in the list preserves these invariants: a
discarded point lies on a line with two chosen points of $S$, so it is not in
$S$ (arc property). At a leaf, $S$ is the chosen set, and it is complete.

The search tree has $816$, $497$, $138$ and $4$ nodes, with $28$, $6$, $22$ and
$0$ leaves. Of these leaves, $7$, $5$, $1$ and $0$ are complete.
\end{proof}

\begin{table}[ht]\centering\small
\begin{tabular}{@{}rll@{}}
\toprule
$j$ & arc (point numbers; each starts with $\mathcal F=0,121,132,12$) & \\
\midrule
1 & 25, 35, 49, 74, 103, 119 & $=R$\\
1 & 25, 35, 53, 59, 94, 109 &\\
1 & 25, 41, 54, 73, 86, 94 &\\
1 & 25, 41, 70, 95, 109, 119 &\\
1 & 25, 46, 70, 85, 93, 117 &\\
1 & 25, 52, 71, 83, 109, 112 &\\
1 & 25, 53, 65, 68, 85, 104 &\\
2 & 26, 41, 49, 69, 90, 105 &\\
2 & 26, 42, 50, 69, 79, 117 &\\
2 & 26, 42, 82, 98, 101, 113 &\\
2 & 26, 43, 62, 79, 91, 104 &\\
2 & 26, 49, 63, 76, 79, 97 &\\
3 & 27, 57, 74, 86, 94, 106 &\\
\bottomrule
\end{tabular}
\caption{The $13$ complete $10$-arcs found by the search of
Lemma~\ref{lem:search}. Each is mapped onto $R$ by an explicit invertible matrix
(\texttt{orbitData} in the Lean file).}\label{tab:arcs}
\end{table}

\begin{lemma}[orbit certificates]\label{lem:cert}
For each arc $T$ of Table~\ref{tab:arcs}, an explicit pair of matrices $(W,W')$
with $WW'=W'W=I$ satisfies $WT=R$.
\end{lemma}
The matrices are listed in the Lean file and were checked there.

\begin{proof}[Proof of Theorem~\ref{thm:main}]
$R$ is a complete $10$-arc (checked directly). Let $A$ be any complete
$10$-arc.
\begin{enumerate}
\item By Lemmas~\ref{lem:frame} and~\ref{lem:preserve}, some $M$ maps $A$ to a
complete $10$-arc through $\mathcal F$.
\item By Lemmas~\ref{lem:reduce}, \ref{lem:search} and~\ref{lem:cert}, a
further projectivity maps that arc onto $R$.
\end{enumerate}
So every complete $10$-arc is equivalent to $R$. As equivalence is symmetric
and transitive (inverse and product matrices), any two complete $10$-arcs are
equivalent.
\end{proof}

\section{Further facts and the literature}
The Python script \texttt{verify.py} computes the following facts. They are not
needed for the proof.
\begin{itemize}
\item \textbf{Arcs through the frame.} For sizes $4,5,\dots,12$, the numbers
of arcs containing $\mathcal F$ are $1$, $72$, $1548$, $10232$, $16233$,
$4032$, $336$, $72$ and $9$.
\item \textbf{Complete ones.} The complete arcs through $\mathcal F$ number
$40, 5103, 3024, 84, 9$ for sizes $7, 8, 9, 10, 12$. So the complete arcs of
$\PG$ have sizes $7,8,9,10,12$.
\item \textbf{Incomplete $10$-arcs.} There are $84$ complete and $252$
incomplete $10$-arcs through $\mathcal F$. Each incomplete one lies on one of
the $9$ conics through $\mathcal F$, and $9\binom{8}{6}=252$. For example,
$C_{10}=\{0,121,132,12,25,42,62,68,81,93\}$ is a $10$-arc whose completion to a
conic adds the points $109$ and $116$. The Lean theorem
\texttt{C10\_not\_complete} proves that $C_{10}$ is not complete.
\item \textbf{One orbit.} The $84$ complete $10$-arcs through $\mathcal F$ have
the same canonical form (least image over all $5040$ frame maps).
\item \textbf{Stabilizer.} The stabilizer of a complete $10$-arc has order $60$,
with element orders $1^1 2^{15} 3^{20} 5^{24}$, those of $A_5$. By orbit
counting, its orbit contains $5040/60=84$ frame-containing arcs, i.e.\ all of
them.
\item \textbf{Total count.} The number of complete $10$-arcs in $\PG$ is
$|\PGL|/60=212427600/60=3540460$.
\end{itemize}

This agrees with the literature. Ball and Lavrauw \cite{bl} (Corollary~8, based
on Hirschfeld's classification \cite[Table~9.4]{hir}) state that for $q=11$
there is a unique complete arc of size $q-1=10$. Their Example~3 is
\[
\{(1,0,0),(0,1,0),(0,0,1),(1,8,7),(1,5,2),(1,3,8),(1,4,1),(1,7,6),(1,9,4),(1,10,5)\}.
\]
\texttt{verify.py} checks that this set is a complete $10$-arc with the same
canonical form as ours. By Segre's theorem \cite{segre} the $12$-arcs are the
conics, and the enumeration shows that each of the $72$ $11$-arcs through
$\mathcal F$ lies on one of the $9$ conics through $\mathcal F$ ($72=9\cdot8$).
So a $10$-arc contained in a larger arc is a subset of a conic, consistent
with the $252$ incomplete ones above.

\section{Lean formalization}
The project \texttt{lean4/} (Lean 4.19.0, core library only) builds in about
$85$ seconds.

\paragraph{Definitions.}
\begin{itemize}
\item \textbf{Linear algebra.} \texttt{F = Fin 11}, vectors \texttt{V},
matrices \texttt{Mat}, \texttt{mulVec}, \texttt{mmul}, \texttt{dot}; and
\texttt{IsInv M N} (two-sided inverse).
\item \textbf{Points and incidence.} \texttt{coord} (the $133$ points),
\texttt{inc}, and \texttt{Collinear} (a common line).
\item \textbf{Arcs.} \texttt{IsArc}, \texttt{IsComplete} and
\texttt{IsCompleteArc k A}, where $A$ is a list of $k$ distinct points.
\item \textbf{Equivalence.} \texttt{act}, \texttt{MapsOnto M A B} (meaning
$\{Ma : a\in A\}=B$), and \texttt{ProjEquiv}.
\end{itemize}

\paragraph{Main theorem.}
\texttt{conjecture\_00000001091} states
\[
(\exists A,\ \mathrm{IsCompleteArc}\ 10\ A)\ \wedge\ \forall A\,B,\
\mathrm{IsCompleteArc}\ 10\ A\to\mathrm{IsCompleteArc}\ 10\ B\to
\mathrm{ProjEquiv}\ A\ B.
\]

\paragraph{General lemmas} (for all matrices and points, no enumeration).
\begin{itemize}
\item \textbf{Algebra.} \texttt{mulVec\_mmul}, \texttt{mmul\_assoc},
\texttt{dot\_mulVec} and \texttt{tr\_mmul}. These are proved from field facts
of $\F$, each checked by \texttt{decide}.
\item \textbf{Action.} \texttt{act\_comp} and \texttt{act\_inv}.
\item \textbf{Collinearity and arcs.} \texttt{collinear\_act},
\texttt{collinear\_of\_act} and \texttt{map\_complete}.
\item \textbf{Search.} The soundness theorem \texttt{search\_sound} and
\texttt{search\_root}, with the counting lemmas.
\item \textbf{Assembly.} \texttt{reduce}, \texttt{to\_rep},
\texttt{frame\_lemma} and \texttt{complete10\_equiv}.
\end{itemize}

\paragraph{Kernel-checked finite facts} (\texttt{decide +kernel}).
\begin{itemize}
\item \texttt{ptl\_check}: the line masks agree with \texttt{inc} on all
$133^2$ pairs (with \texttt{inc\_iff}, proved algebraically).
\item \texttt{T1ok}--\texttt{T4ok}: the frame-lemma tables.
\item \texttt{g1ok}--\texttt{g4ok}: the frame stabilizer data.
\item \texttt{cand\_cover}, \texttt{cover1}--\texttt{cover4} and
\texttt{rest\_ok}.
\item \texttt{search1}--\texttt{search4}: the four searches.
\item \texttt{orbit\_ok}: the $13$ certificates.
\item \texttt{R\_complete}.
\end{itemize}

\paragraph{Model sanity.} The theorems \texttt{points\_cover},
\texttt{points\_distinct} and \texttt{coord\_ne\_zero} show that the $133$
listed points are exactly the points of $\PG$. \texttt{C10\_not\_complete}
shows that completeness is a genuine restriction.

\paragraph{Axioms.} \texttt{\#print axioms} reports only \texttt{propext} and
\texttt{Quot.sound}. There is no \texttt{sorry}, no \texttt{native\_decide} and
no added axiom.

\paragraph{Limitations.} Collinearity is defined as ``lying on a common line''
(the incidence definition of $\PG$). The equivalence with the vanishing of the
$3\times3$ determinant is standard, but it is not proved in Lean. The facts of
Section~4 (the count $84$, the stabilizer $A_5$, the size spectrum) are
checked only by \texttt{verify.py}.

\section{Independent check}
\texttt{verify.py} uses only the Python standard library and no Lean table. It
uses a different normalization (last nonzero coordinate $1$) and decides
collinearity by the determinant. It:
\begin{enumerate}
\item enumerates all arcs through the frame;
\item re-checks the $84$ complete $10$-arcs, and that each of the $252$
incomplete ones lies on a conic;
\item computes the canonical forms ($1$ class) and the stabilizer (order $60$,
$A_5$);
\item checks the Ball--Lavrauw example;
\item reproduces the statistics of the Lean search ($1455$ nodes, $56$ leaves,
$13$ complete).
\end{enumerate}
It runs in about $20$ seconds.

\section{Reproduction}
\begin{verbatim}
cd lean4 && lake build        # Lean 4.19.0, core only
cd .. && python3 verify.py
pdflatex report.tex && pdflatex report.tex
\end{verbatim}

\begin{thebibliography}{9}
\bibitem{bl} S.~Ball and M.~Lavrauw, Planar arcs, \emph{J. Combin. Theory
Ser. A} (2018); arXiv:1705.10940.
\bibitem{hir} J.~W.~P.~Hirschfeld, \emph{Projective Geometries over Finite
Fields}, 2nd ed., Oxford University Press, 1998.
\bibitem{segre} B.~Segre, Ovals in a finite projective plane, \emph{Canad. J.
Math.} 7 (1955), 414--416.
\end{thebibliography}
\end{document}
